% !TeX root = ../../Thesis.tex
\ThesisChapter
  {\ThesisText{Concepts and Architecture}{Konzepte und Architektur}}
  {ch:concepts}
  {TODO: Add a brief chapter overview. \\
  This chapter derives requirements and describes the concepts, design decisions, and resulting architecture.}

% TODO: Describe the concepts and architecture developed or used in this thesis.

This chapter presents the concrete solution approach of the thesis.
It should explain the main concepts, design decisions, components, interfaces, and interactions that form the proposed system or method.

Possible contents include:
\begin{itemize}
  \item the requirements and design goals,
  \item the main concepts of the proposed solution,
  \item the overall system architecture,
  \item the responsibilities of individual components,
  \item the interfaces and communication between components,
  \item relevant data flows or processing steps, and
  \item important architectural and design decisions.
\end{itemize}

Clearly explain how the proposed concepts address the problem introduced earlier in the thesis.
Diagrams may be used to illustrate the architecture, components, interfaces, and data flows.

This chapter should focus on the specific solution developed or applied in the thesis.
General theories, technologies, terminology, and other knowledge required to understand the work belong in the background chapter.
Implementation-specific details should be presented in the implementation chapter.

% Example for a diagram drawn directly in LaTeX with TikZ.
\section{\ThesisText{Example Architecture Diagram}{Beispiel: Architekturdiagramm}}
\label{sec:example-architecture-diagram}

For diagrams and other freely designed graphics, draw.io / diagrams.net may be used:~\url{https://app.diagrams.net/}

\textbf{Important:} Every figure generated entirely or partially using an AI-based tool must be explicitly identified and properly referenced.
The use of AI must be clearly disclosed in the corresponding figure caption.

Figure~\ref{fig:example-architecture} sketches a scheduler that distributes work to several compute nodes.
Diagrams of this size can be drawn directly in LaTeX with TikZ, which keeps them vector-based and matches the fonts and colors of the surrounding text.
The source is stored in the chapter-specific figure folder and included like any other asset:

% !TeX root = ../../Thesis.tex
\begin{figure}[htbp]
  \centering
  \begin{tikzpicture}[
    every node/.style={font=\small},
    unit/.style={
      draw, rounded corners=2pt, line width=.7pt,
      minimum width=20mm, minimum height=8mm
    },
    master/.style={unit, draw=hpcgreen, fill=hpcgreen!8},
    worker/.style={unit, draw=hpcblue, fill=hpcblue!8},
    wire/.style={draw=hpcgray, line width=.7pt},
    feed/.style={wire, ->, >=stealth}
  ]
    \node[master] (sched) {Scheduler};
    \node[worker, below=11mm of sched] (n2) {Node 2};
    \node[worker, left=7mm of n2] (n1) {Node 1};
    \node[worker, right=7mm of n2] (n3) {Node 3};

    % Vertical trunk, horizontal bus, then one drop per node.
    \coordinate (bus) at ($(sched.south)!.5!(n2.north)$);
    \draw[wire] (sched.south) -- (bus);
    \draw[wire] (n1.north |- bus) -- (n3.north |- bus);
    \foreach \n in {n1,n2,n3}{
      \draw[feed] (\n.north |- bus) -- (\n.north);
    }
  \end{tikzpicture}
  \caption{\ThesisText{Example architecture diagram drawn with TikZ: a scheduler distributing work to three compute nodes.}{Beispielhaftes Architekturdiagramm in TikZ: ein Scheduler verteilt Arbeit auf drei Rechenknoten.}}
  \label{fig:example-architecture}
\end{figure}

